% answers_short.tex —— 2012 年·天津卷·理科 · 简短答案速查
% 由 tools/make_short_answers.py 自 answers.tex 的 \gksolhead/\gkanswer 提取，
% 并与 2012/2012_数学真题_天津理_答案要点.md 逐题比对后生成，勿手改（改源 MD 后重跑）。

\begin{gknotice}
\textbf{使用说明}：本节只给最终答案便于快速核对，完整推导见解析版。答案为编写者独立推导，\textbf{非官方标准答案}。
\end{gknotice}

\gkpart{选择题}
\begin{center}\begin{tabular}{ccccccccc}
\toprule
题号 & 1 & 2 & 3 & 4 & 5 & 6 & 7 & 8 \\
\midrule
答案 & B & A & C & B & D & A & A & C \\
\bottomrule
\end{tabular}\end{center}

\gkpart{填空题}
\begin{itemize}[leftmargin=2.6em]
  \item[\textbf{9.}] $18$；$9$。
  \item[\textbf{10.}] $18+9\pi$。
  \item[\textbf{11.}] $m=-1$；$n=1$。
  \item[\textbf{12.}] $2$。
  \item[\textbf{13.}] $\dfrac43$。
  \item[\textbf{14.}] $(0,1)\cup(1,4)$。
\end{itemize}

\gkpart{解答题}
\begin{itemize}[leftmargin=2.6em]
  \item[\textbf{15.}] (1) $T=\pi$；(2) 最大值 $\sqrt2$，最小值 $-1$
  \item[\textbf{16.}] (1) $\dfrac8{27}$；(2) $\dfrac19$；(3) $P(\xi=0)=\dfrac8{27}$，$P(\xi=2)=\dfrac{40}{81}$，$P(\xi=4)=\dfrac{17}{81}$；$E\xi=\dfrac{148}{81}$
  \item[\textbf{17.}] (1) 证明见解析；(2) $\dfrac{\sqrt{30}}6$；(3) $\dfrac{\sqrt{10}}{10}$
  \item[\textbf{18.}] (1) $a_n=3n-1$，$b_n=2^n$；(2) 证明见解析
  \item[\textbf{19.}] (1) $\dfrac{\sqrt2}2$；(2) 证明见解析
  \item[\textbf{20.}] (1) $a=1$；(2) $\dfrac12$；(3) 证明见解析
\end{itemize}
